\subsection{How long the Collatz probes can take}

The bracket is observable, but its hitting time necessarily depends on
projective mixing of the current face.

\begin{proposition}[Projective Collatz hitting time]
\label{prop:collatz-projective-hitting}
Let $M>0$ be the normalized residual map in
\cref{thm:collatz-gated-window}, let $\phi>0$ be its unit Perron vector, and
write $d_H$ for Hilbert's projective metric.  Define
\begin{equation*}
 \Delta(M):=\log\max_{i,j,k,l}
 \frac{M_{ik}M_{jl}}{M_{il}M_{jk}},
 \qquad
 \tau(M):=\tanh\!\left(\frac{\Delta(M)}4\right)<1.
\end{equation*}
For $z_s=M^sz_0$,
\begin{equation}
 \log\frac{U_s}{L_s}
 \leq(1+\tau)\tau^s d_H(z_0,\phi).
 \label{eq:collatz-hilbert-width}
\end{equation}
Consequently the bracket condition in
\cref{thm:collatz-gated-window} is reached after
\begin{equation}
 O\!\left(
 \frac{\log(d_H(z_0,\phi)/q_{\rm face}^2)}{-\log\tau(M)}
 \right)
 \label{eq:collatz-hilbert-hitting}
\end{equation}
probes.

Alternatively, let $\delta=\mu_2/\mu_1<1$ and
$\gamma=\min_i\phi_i$.  For every $z_0\geq0$, $z_0\neq0$, the observable
stop is reached by
\begin{equation}
 P\leq
 \left\lceil
 \frac{\log(128/(\gamma^2q_{\rm face}^2))}{-\log\delta}
 \right\rceil.
 \label{eq:collatz-spectral-hitting}
\end{equation}
The algorithm need not know $\delta$ or $\gamma$; they occur only in this
proof-only work bound.
\end{proposition}

\begin{proof}
Birkhoff contraction gives
$d_H(z_s,\phi)\leq\tau^sd_H(z_0,\phi)$.  Since
\begin{equation*}
 \log(U_s/L_s)=d_H(Mz_s,z_s),
\end{equation*}
the triangle inequality and one further contraction prove
\cref{eq:collatz-hilbert-width}.  Once its right-hand side is at most
$q_{\rm face}^2/128$, elementary exponential bounds and
$L_s\leq\mu_1\leq U_s$ imply
$U_s-L_s\leq(1-U_s)/16$.

For the second statement, write
$z_0=c_1\phi+z_\perp$.  Positivity gives
$c_1=\langle\phi,z_0\rangle\geq\gamma\norm{z_0}_2$, while
\begin{equation*}
 \norm{\mu_1^{-s}M^sz_\perp}_2
 \leq\delta^s\norm{z_0}_2.
\end{equation*}
Thus every coordinate has relative Perron error at most
$\delta^s/\gamma^2$.  Making this at most
$q_{\rm face}^2/128$ proves \cref{eq:collatz-spectral-hitting} and the same
bracket stop.
\end{proof}

There is no graph-uniform $\widetilde O(1/q_{\rm face})$ bound for this probe
phase.  On full paths one can have
$q_{\rm face}=\Theta(n^{-1})$,
$\delta=1-\Theta(n^{-2})$, and $\gamma^2=\Theta(n^{-1})$.
A positive start with a nonzero second-mode component then needs
$\Omega(n^2\log n)$ pure-prox probes before the bracket has width
$\Theta(q_{\rm face}^2)$.  Explicitly scanning and applying the face map
costs $\Omega(n)$ per probe.  The observable stop remains a correct
conditional accelerator, but without a face mixing certificate its
pre-window work can be $\Omega(n^3\log n)$.
